El presente trabajo de grado desarrolló e implementó \textbf{ViaData — Medellín}, instrumento virtual geoespacial orientado al análisis y la proyección exploratoria de la accidentalidad vial en Medellín a partir de datos abiertos Mede. A continuación se sintetizan las conclusiones derivadas del diseño, la implementación, la evaluación numérica y la validación técnica del sistema, en correspondencia con el objetivo general y los objetivos específicos planteados.

En relación con el \textbf{objetivo general}, se logró construir una plataforma web que integra visualización interactiva, segmentación territorial y temporal, e indicadores y modelos estadísticos transparentes para apoyar la lectura del fenómeno y la discusión de escenarios en seguridad vial. El tablero y el mapa permiten consulta pública de KPIs, series, matrices día$\times$hora y capas geoespaciales; el módulo de predicciones, los reportes y el asistente conversacional amplían esa capacidad para perfiles analista y administrador, con trazabilidad entre la interfaz y el motor de cálculo del backend.

Respecto al \textbf{objetivo específico 1} (arquitectura y modelo de datos geoespacial), se adoptó una arquitectura monolítica en Django con API REST, autenticación JWT y frontend React, apoyada en PostgreSQL/PostGIS. El flujo ETL reproducible ---centralizado en \texttt{mede\_limpieza.py} y documentado en manuales de carga--- transforma el Excel Mede en una base relacional con incidentes, víctimas, catálogos y geometrías, habilitando consultas espaciales y el contraste entre modo territorio \textit{registro} y \textit{espacial}.

En cuanto al \textbf{objetivo específico 2} (módulos de análisis estadístico), el sistema calcula indicadores descriptivos y diagnósticos (densidad G01--G02, calidad G03, rankings, matrices temporales) y un módulo predictivo estructurado en cinco bloques (§1--§5). La evaluación cerrada y el caso de estudio escenario A adoptaron: \textbf{SARIMA}$(2,1,3)(1,1,1)_{12}$ para proyección mensual de incidentes a nivel ciudad (MAPE hold-out 12,62\,\%); \textbf{índice P05} para priorización territorial del periodo; \textbf{modelo estacional} para ranking de carga futura por comuna; \textbf{logit con exposición} para la proporción de víctimas fatales (MAPE hold-out 20,25\,\%); y \textbf{reparto histórico} con suavizado Laplace para patrones día$\times$hora y día de la semana, heredando el total de §1.

Frente al \textbf{objetivo específico 3} (tablero de monitoreo y diagnóstico), la ruta \texttt{/tablero} materializa la evolución mensual, la comparación interanual y la distribución por territorio, día y hora. El mapa analítico complementa esa lectura con puntos, densidad, coroplética y concentraciones espaciales. Los resultados confirman concentraciones recurrentes ---La Candelaria en volumen e índice P05; martes 07:00 en patrones temporales--- que orientan la focalización territorial.

En el \textbf{objetivo específico 4} (validación técnica), la suite \texttt{pytest} sobre las aplicaciones críticas del backend finalizó en verde, y el módulo de predicciones cuenta con evaluación reproducible en CSV y scripts documentados. El criterio de validación predictiva basado en \textit{hold-out} (MAPE $\leq 20\,\%$ como umbral exploratorio) resultó más informativo que exigir R$^2 \geq 0{,}80$ en series mensuales con estacionalidad y confinamiento COVID.

Además de los objetivos formales, se entregaron componentes complementarios ---mapa analítico, reportes imprimibles, asistente con Gemini y gestión de usuarios--- que facilitan la operación y la comunicación de hallazgos sin alterar la formulación académica del proyecto.

Desde el plano analítico, las conclusiones centrales del estudio son las siguientes:

\begin{enumerate}
	\item Los datos abiertos Mede, depurados y cargados con reglas únicas, permiten sostener un instrumento geoespacial consultable sobre del orden de 160\,000+ incidentes georreferenciados en el periodo cargado (aprox. 2014--2021).
	\item La analítica descriptiva y geoespacial evidencia que la siniestralidad se concentra en corredores y comunas del valle urbano, con patrones temporales persistentes que complementan la lectura por volumen bruto.
	\item \textbf{No existe un modelo predictivo único óptimo} para todos los escenarios: la efectividad depende de la variable, el territorio, la longitud del periodo y el tratamiento del shock COVID; la interfaz expone esa contingencia mediante comparación multi-modelo bajo filtros explícitos.
	\item El hold-out discrimina mejor que el ajuste in-sample entre modelos parsimoniosos y especificaciones más flexibles; Poisson y estacional pueden ajustar bien el historial y perder en meses reservados, mientras SARIMA lidera en incidentes ciudad bajo escenario A.
	\item Las proyecciones son \textbf{herramientas exploratorias} para priorización y planificación: apoyan la discusión de escenarios bajo estabilidad del patrón histórico filtrado, pero no establecen causalidad ni estiman riesgo individual \cite{oms2024,parraga2024}.
\end{enumerate}

En conjunto, ViaData — Medellín demuestra que es viable articular ingeniería de software, sistemas de información geográfica y ciencia de datos en un mismo entorno orientado a la seguridad vial urbana, con modelos interpretables y evaluación documentada \cite{lugo2025,monar2025}. Las limitaciones de generalización temporal, precisión territorial fina, ausencia de variables exógenas y alcance omitido (ranking por vía, ingesta automática, ML espacial avanzado) se abordaron en la discusión y orientan las líneas de trabajo futuro de la siguiente sección.
